%! Graphing Quadratic Functions & Transformations — Unit 2, Lesson 2.1 — Warm-Up (Answer Key)
% Mirrors ../warmup/main.tex EXACTLY (12pt, ONE page): every \blank{} becomes an \ans{} and
% both work blocks are byte-identical.
% The warm-up is set at 12pt like every other student component, and it must still fit ONE page,
% blank and keyed. Three spiral-review items rehearsing exactly what today leans on: item 1 is
% FOIL, and its two expansions quietly plant the hook (both give x^2-2x-3); item 2 brings
% Lesson 1.3's transformation language of y=|x| into the room before the parabola gets it;
% item 3 is the hand-off, LEFT HANGING — notes row 1 answers it.
% Both work blocks are BYTE-IDENTICAL in warmup_key/main.tex.
\documentclass[12pt]{article}
\usepackage{algebra2-article}
\usepackage{algebra2-key}   % pulls in -boxes; do NOT also load -boxes
% Coordinate grid: {xmin}{xmax}{ymin}{ymax}
\newcommand{\qgrid}[4]{%
  \draw[step=1, linegray!40, very thin] (#1,#3) grid (#2,#4);
  \draw[->, semithick, charcoal] (#1,0)--(#2,0) node[below right, font=\scriptsize]{$x$};
  \draw[->, semithick, charcoal] (0,#3)--(0,#4) node[left, font=\scriptsize]{$y$};}
\newcommand{\fdot}[3]{\fill[#1] (#2,#3) circle (3.0pt);}

\begin{document}

\pageheader{Unit 2, Lesson 2.1}{Warm-Up --- Answer Key}
% NAMESTRIP: no \namedateperiod here — mirrors the blank.

\vspace{0.08in}

\begin{notesbox}{Getting Ready: Skills We'll Use Today}
\small
Three things you already know. Today one parabola will show up wearing three different
outfits --- work quickly, then check with a neighbour.

\vspace{5pt}
\textbf{1. Multiply it out.} Expand each expression and simplify.

\begin{minipage}[t]{0.47\linewidth}
\begin{work}
  (x-1)^2-4 &= (x-1)(x-1)-4 \\
            &= x^2-2x+1-4 \\
            &= x^2-2x-3
\end{work}
\end{minipage}\hfill
\begin{minipage}[t]{0.47\linewidth}
\begin{work}
  (x+1)(x-3) &= x^2-3x+x-3 \\
             &= x^2-2x-3
\end{work}
\end{minipage}

\vspace{2pt}
Did the two expressions give the \emph{same} answer? \ans{yes} \quad So how many different
curves are hiding in those two lines? \ans{one}

\vspace{7pt}
\textbf{2. Moving a parent graph (Lesson 1.3).} The parent is $y=|x|$, the V with its corner at
the origin.
\begin{itemize}[leftmargin=1.6em, itemsep=3pt]
  \item $y=|x|+2$ moves the parent \ans{up} (up / down) by \ans{$2$}.
  \item $y=|x-4|$ moves the parent \ans{right} (left / right) by \ans{$4$}.
  \item $y=-|x|$ opens \ans{down} (up / down), so its corner is a \ans{maximum}
        (maximum / minimum).
\end{itemize}

\vspace{7pt}
\textbf{3. Hold on to this one.} Below are the parent $y=x^2$ and a second graph, $y=(x-3)^2$.
\begin{center}
\begin{tikzpicture}[scale=0.34]
  \qgrid{-3.5}{3.5}{-0.5}{8.5}
  \begin{scope}
    \clip (-3.5,-0.5) rectangle (3.5,8.5);
    \draw[very thick, charcoal, domain=-3:3, samples=70, smooth, variable=\x] plot (\x,{(\x)^2});
  \end{scope}
  \fdot{charcoal}{0}{0}
  \node[font=\scriptsize\bfseries, charcoal] at (0,9.6) {$y=x^2$};
  \begin{scope}[xshift=10cm]
    \qgrid{-1.5}{6.5}{-0.5}{8.5}
    \begin{scope}
      \clip (-1.5,-0.5) rectangle (6.5,8.5);
      \draw[very thick, forest, domain=0:6, samples=70, smooth, variable=\x] plot (\x,{(\x-3)^2});
    \end{scope}
    \fdot{forest}{3}{0}
    \node[font=\scriptsize\bfseries, forest] at (2.5,9.6) {$y=(x-3)^2$};
  \end{scope}
\end{tikzpicture}
\end{center}
\vspace{-2pt}
The second graph moved \ans{right} (left / right) by \ans{$3$}. \quad Does that surprise
you, given the minus sign? \ans{yes --- minus inside moves it \emph{right}}
\end{notesbox}

\end{document}
